\documentclass[../main.tex]{subfiles}
\begin{document}
    \begin{enumerate}[a), font=\bfseries]
        \item Calculer l'aire du rectangle \textsc{ABCD} de \textbf{deux manières différentes}.
        \begin{center}
            \begin{tikzpicture}[scale=0.8]
                \def\L{6}
                \def\l{3}
                \def\ll{2}
                \coordinate[label=below left:{A}](A)at(0,0);
                \coordinate[label=below right:{B}](B)at(\L,0);
                \coordinate[label=above right:{C}](C)at(\L,\l);
                \coordinate[label=above left:{D}](D)at(0,\l);
                \coordinate[label=below:{E}](E)at(\ll,0);
                \coordinate[label=above:{F}](F)at(\ll,\l);
                \tkzDrawPolygon(A,B,C,D)
                \draw(E)--(F);
                \path(A)--(D)node[midway, sloped, above]{3 cm};
                \path(D)--(F)node[midway, sloped, above]{2 cm};
                \path(F)--(C)node[midway, sloped, above]{4 cm};
            \end{tikzpicture}
        \end{center}
        \item Écrire \textbf{deux expression littérale} donnant l'aire du rectangle \textsc{ABCD}.
        \begin{center}
            \begin{tikzpicture}[scale=0.8]
                \def\L{6}
                \def\l{3}
                \def\ll{2}
                \coordinate[label=below left:{A}](A)at(0,0);
                \coordinate[label=below right:{B}](B)at(\L,0);
                \coordinate[label=above right:{C}](C)at(\L,\l);
                \coordinate[label=above left:{D}](D)at(0,\l);
                \coordinate[label=below:{E}](E)at(\ll,0);
                \coordinate[label=above:{F}](F)at(\ll,\l);
                \tkzDrawPolygon(A,B,C,D)
                \draw(E)--(F);
                \path(A)--(D)node[midway, left]{k};
                \path(D)--(F)node[midway, sloped, above]{a};
                \path(F)--(C)node[midway, sloped, above]{b};
            \end{tikzpicture}
        \end{center}
        \item En vous aidant de la question précédente, compléter la propriété ci dessous:~\\
        \begin{simplebox}
            Pour tout nombre $a$, $b$, $k$: $$k(a+b)=$$
            \vspace{.3cm}
        \end{simplebox}
        \item \textbf{Application :} développer les expression littérales suivantes:
        \begin{center}
            \begin{minipage}{.45\linewidth}
                $3(x+4)=$
            \end{minipage}\hspace{.05\linewidth}
            \begin{minipage}{.45\linewidth}
                $-2\times(1+x)=$
            \end{minipage}
        \end{center}
    \end{enumerate}
\end{document}